% Una subsección por cada problema/algoritmo analizado.
Presente brevemente qué resultados se muestran en esta sección.

\subsection{Primer experimento}

Describa los resultados y refiérase a las figuras con \verb|\cref{fig:ejemplo}|.

% --- Figura simple ---
% \begin{figure}[H]
%   \centering
%   \includegraphics[width=0.7\textwidth]{nombre_imagen.png}
%   \caption{Descripción de la figura.}
%   \label{fig:ejemplo}
% \end{figure}

% --- Dos figuras lado a lado ---
% \begin{figure}[H]
%   \centering
%   \begin{minipage}[t]{0.5\textwidth}
%     \includegraphics[width=\textwidth]{grafico_lineas.png}
%   \end{minipage}%
%   \begin{minipage}[t]{0.5\textwidth}
%     \includegraphics[width=\textwidth]{grafico_puntos.png}
%   \end{minipage}
%   \caption{Descripción de las figuras.}
%   \label{fig:ejemplo2}
% \end{figure}

% --- Tabla ---
% \begin{table}[H]
%   \centering
%   \begin{tabular}{lrr}
%     \toprule
%     Algoritmo & Tiempo (s) & Memoria (KB) \\
%     \midrule
%     A & 0.01 & 1000 \\
%     B & 0.02 & 2000 \\
%     \bottomrule
%   \end{tabular}
%   \caption{Descripción de la tabla.}
%   \label{tab:ejemplo}
% \end{table}

\subsection{Segundo experimento}

Análisis y comparación con lo esperado teóricamente.
